% ================= SEKCJA 3 =================
\section{Cel projektu}

Celem projektu \texttt{WorkFlow} jest zaprojektowanie oraz wykonanie aplikacji webowej wspierającej zarządzanie pracownikami, czasem pracy, dokumentacją kadrową oraz danymi wykorzystywanymi w procesie rozliczeń. System ma uporządkować informacje, które w wielu organizacjach są przechowywane w rozproszony sposób, na przykład w arkuszach kalkulacyjnych, dokumentach papierowych, wiadomościach e-mail lub niezależnych plikach.

Projekt zakłada utworzenie jednego środowiska, w którym możliwe jest prowadzenie ewidencji pracowników, zakładów, projektów, umów, certyfikatów, dokumentów, nieobecności oraz wpisów czasu pracy. Dane te są ze sobą powiązane, dlatego ważnym celem systemu jest zachowanie spójności pomiędzy poszczególnymi obszarami aplikacji.

\subsection{Uporządkowanie danych kadrowych}

Jednym z podstawowych celów projektu jest stworzenie centralnej ewidencji pracowników. Każdy pracownik posiada profil zawierający dane identyfikacyjne, przypisanie do zakładu, rolę systemową, status aktywności oraz opcjonalny identyfikator karty RFID wykorzystywany przy rejestracji czasu pracy.

Profil pracownika stanowi punkt wyjścia dla pozostałych modułów systemu. Z tego poziomu można powiązać pracownika z umowami, certyfikatami, dokumentami, nieobecnościami, projektami oraz wpisami czasu pracy. Takie podejście ogranicza powielanie danych i ułatwia odnalezienie informacji związanych z daną osobą.

\subsection{Obsługa struktury organizacyjnej}

System powinien umożliwiać odwzorowanie podstawowej struktury organizacyjnej firmy. W projekcie przyjęto podział na zakłady, do których przypisywani są pracownicy, projekty oraz czytniki RCP. Dodatkowo przewidziano mechanizm dostępu użytkowników do wybranych zakładów.

Celem tego rozwiązania jest umożliwienie pracy w organizacji posiadającej więcej niż jedną jednostkę organizacyjną. Dzięki temu użytkownik może mieć dostęp tylko do danych tych zakładów, za które odpowiada.

\subsection{Rejestracja i weryfikacja czasu pracy}

Istotnym celem projektu jest obsługa czasu pracy. System umożliwia przyjmowanie surowych zdarzeń z czytników RCP, takich jak wejście pracownika do pracy oraz wyjście z pracy. Zdarzenia te są zapisywane w bazie jako rekordy typu \texttt{TimeEvent}.

Na podstawie pary zdarzeń \texttt{IN} oraz \texttt{OUT} system tworzy wpis czasu pracy. Wpis ten zawiera godzinę rozpoczęcia, godzinę zakończenia, wyliczoną liczbę godzin, pracownika oraz projekt, z którym powiązany jest czytnik. Dodatkowo wpis czasu pracy posiada status, dzięki czemu można odróżnić dane oczekujące na weryfikację od danych zatwierdzonych.

\subsection{Obsługa dokumentacji pracowniczej}

Kolejnym celem projektu jest uporządkowanie dokumentacji związanej z pracownikami. System przewiduje obsługę dokumentów, umów, certyfikatów, badań, szkoleń oraz nieobecności.

Dokumenty są przechowywane jako pliki w magazynie obiektowym, natomiast baza danych przechowuje ich metadane oraz powiązania z pracownikiem, umową lub certyfikatem. Takie rozwiązanie oddziela przechowywanie plików od danych relacyjnych, a jednocześnie pozwala zachować logiczne powiązania pomiędzy dokumentami a rekordami systemu.

\subsection{Kontrola ważności certyfikatów i badań}

System ma wspierać kontrolę terminów ważności certyfikatów, szkoleń, badań lekarskich oraz innych uprawnień pracowników. W tym celu w projekcie przewidziano słownik certyfikatów oraz tabelę konkretnych certyfikatów przypisanych do pracowników.

Celem tego modułu jest umożliwienie szybkiego sprawdzenia, które dokumenty są aktualne, które zbliżają się do końca ważności oraz które są już przeterminowane. Funkcja ta jest szczególnie istotna w przypadku dokumentów wymaganych do wykonywania określonych zadań lub pracy na konkretnych stanowiskach.

\subsection{Rejestracja nieobecności}

Projekt obejmuje również obsługę nieobecności pracowników. System pozwala zapisać typ nieobecności, datę rozpoczęcia, datę zakończenia, status zatwierdzenia oraz opcjonalny dokument powiązany z nieobecnością.

Uwzględniono kilka podstawowych typów nieobecności, takich jak urlop, zwolnienie lekarskie, nieobecność nieusprawiedliwiona oraz nieobecność specjalna. Dane te uzupełniają profil pracownika i mogą być wykorzystane podczas analizy obecności oraz planowania pracy.

\subsection{Przygotowanie danych do rozliczeń}

Celem systemu nie jest zastąpienie pełnego systemu księgowego, lecz przygotowanie uporządkowanych danych, które mogą zostać wykorzystane przez osoby odpowiedzialne za rozliczenia. Podstawą tych danych są zatwierdzone wpisy czasu pracy oraz informacje o aktualnych umowach pracowników.

System umożliwia utworzenie eksportu rozliczeniowego, który grupuje zatwierdzone wpisy czasu pracy dla danego okresu. Dzięki temu można ograniczyć ryzyko wielokrotnego rozliczenia tych samych godzin oraz zachować informację o tym, kto wygenerował daną paczkę rozliczeniową.

\subsection{Kontrola dostępu i role użytkowników}

W projekcie zastosowano podział użytkowników według ról. System przewiduje między innymi role administratora, pracownika HR, księgowości, brygadzisty, pracownika oraz użytkownika biurowego.

Celem tego mechanizmu jest ograniczenie dostępu do funkcji zgodnie z zakresem odpowiedzialności danej osoby. Administrator posiada najszerszy zakres uprawnień, HR odpowiada za dane kadrowe, księgowość za dane rozliczeniowe, a brygadzista za informacje związane z pracownikami i projektami przypisanymi do jego obszaru działania.

\subsection{Cel techniczny projektu}

Oprócz realizacji wymagań funkcjonalnych projekt ma również cel techniczny. System został przygotowany jako aplikacja działająca w środowisku kontenerowym. W skład środowiska wchodzą: frontend, backend, baza PostgreSQL, magazyn plików MinIO oraz Nginx pełniący rolę reverse proxy.

Takie podejście ułatwia uruchomienie aplikacji w powtarzalny sposób. Projekt zawiera również konfigurację środowiska, przykładowy plik zmiennych oraz skrypty seedujące bazę. Dzięki temu możliwe jest uruchomienie czystej instancji systemu albo instancji zawierającej dane demonstracyjne potrzebne do prezentacji.

\subsection{Zakładany rezultat}

Rezultatem projektu ma być działająca aplikacja webowa, która prezentuje spójny model zarządzania pracownikami, czasem pracy, dokumentacją i rozliczeniami. System powinien umożliwiać przejście przez podstawowy proces: od utworzenia pracownika, przez przypisanie go do zakładu i projektu, rejestrację czasu pracy, obsługę dokumentów i certyfikatów, aż po przygotowanie danych do rozliczeń.

Projekt stanowi podstawę do dalszego rozwoju. W przyszłości może zostać rozszerzony między innymi o bardziej zaawansowane raportowanie, powiadomienia, integrację z zewnętrznym systemem księgowym lub zewnętrznym dostawcą tożsamości.